% Mục 1 (bản rút gọn). Rút từ sections/intro.tex.
Nhóm cài đặt bằng C++ mutex lock (TAS, TTAS, MCS), sense-reversing barrier và concurrent counter
(lock, CAS, fetch-and-add). Nhóm kiểm tra tính đúng đắn của các cài đặt và đo khả năng mở rộng trên máy
shared-memory, so với cấu trúc của thư viện. Phần 4 của đề (MPI-3 RMA~\cite{mpi41}) không bắt buộc và nhóm không làm.

\begin{itemize}[nosep]
  \item \emph{Cache line.} Đơn vị (128 byte trên hai máy Apple, 64 byte trên \PCname{LAP1184})
    mà phần cứng giữ nhất quán giữa các core. Nhiều core cùng đọc được, nhưng ghi thì một core phải giữ
    line độc quyền và làm mất hiệu lực bản sao ở core khác.
  \item \emph{Vòng lặp đo.} Mỗi vòng, thread lấy lock, gọi \texttt{is\_prime(N)} trong critical
    section, nhả lock, rồi gọi \texttt{is\_prime(N)} thêm OW lần ở ngoài (OW viết tắt
    \texttt{OUTSIDE\_WORK}; $\mathrm{OW} = 0$ là tranh chấp tối đa). $p$ là số thread.
  \item \emph{TAS, TTAS, MCS.} TAS lặp \texttt{exchange} lên cờ lock. TTAS đọc cờ tới khi rảnh rồi mới
    \texttt{exchange}~\cite{anderson1990spinlock}. MCS xếp hàng, mỗi thread chờ trên cờ
    riêng~\cite{mellorcrummey1991mcs}.
  \item \emph{Episode.} Một lần cả $p$ thread cùng vượt barrier~\cite{hensgen1988barrier}.
  \item \emph{FAA, CAS.} Fetch-and-add cộng và trả giá trị cũ trong một bước. Compare-and-swap chỉ ghi
    nếu giá trị chưa đổi, nếu không thì thất bại và phải thử lại.
\end{itemize}

Mục Kết quả giải thích số đo chủ yếu bằng số lần một cache line đi giữa các core trên đường tới
hạn~\cite{herlihy2020art}, cùng với scheduler và hai loại core của máy. P-core là core hiệu năng cao.
E-core là core tiết kiệm điện, chậm hơn, và hai loại nằm ở hai cụm có cache L2 riêng. Các hình đánh dấu
số P-core và số core của máy (trên \PCname{apple-m5} là \PCval{apple-m5}{pcores} và
\PCval{apple-m5}{ncpu}).

\begin{description}[leftmargin=0pt, itemsep=0.2em, font=\normalfont]
  \item[RQ1 (đúng đắn).] Các cài đặt có đúng không, và bộ kiểm tra có đủ nhạy để bắt lỗi cài đặt?
  \item[RQ2 (lock: mở rộng).] TAS, TTAS, MCS mở rộng thế nào theo số thread, \texttt{OUTSIDE\_WORK} và
    $N$, so với lock thư viện?
  \item[RQ3 (lock: backoff và công bằng).] Backoff giúp hay hại, và ở chế độ nào? Lock nhanh nhất có
    công bằng không?
  \item[RQ4 (barrier).] Chi phí một episode tăng thế nào theo số thread, theo cách chờ, so với barrier
    thư viện?
  \item[RQ5 (counter).] Thứ tự hiệu năng giữa FAA, CAS (có/không backoff) và counter dùng lock, và vì
    sao?
  \item[RQ6 (giữa các máy).] Các kết luận trên có giữ nguyên trên máy khác không?
\end{description}

\begin{table}[H]\centering\footnotesize
\begin{tabularx}{\linewidth}{@{}l>{\raggedright\arraybackslash}p{4.4cm}>{\raggedright\arraybackslash}X@{}}
\toprule
RQ & Thí nghiệm & Hình, bảng \\
\midrule
RQ1 & stress, TSan, mutation & Bảng~\ref{tab:corr:mutation} \\
RQ2 & lock $\times$ thread $\times$ OW $\times$ $N$, khoảng trống trao lock &
  Hình~\ref{fig:apple-m5:lock_slowdown}; Bảng~\ref{tab:handoff} \\
RQ3 & sáu biến thể backoff & Hình~\ref{fig:apple-m5:lock_backoff},
  \ref{fig:apple-m5:lock_fairness} \\
RQ4 & episode/giây, bảy barrier & Hình~\ref{fig:apple-m5:barrier_rate} \\
RQ5 & thời gian, số CAS hỏng & Hình~\ref{fig:apple-m5:counter_time},
  \ref{fig:apple-m5:counter_cas_failures} \\
RQ6 & cùng quy trình trên \PCname{C4M7Y670NT} và \PCname{LAP1184} & Hình~\ref{fig:cross:lock},
  \ref{fig:cross:barrier}, \ref{fig:cross:counter}; Bảng~\ref{tab:cross:machines} \\
\bottomrule
\end{tabularx}
\caption{Câu hỏi nghiên cứu, thí nghiệm và hình, bảng trả lời.}\label{tab:intro:rqmap}
\end{table}

Mục Nền tảng thí nghiệm mô tả máy đo, và mục Cài đặt trình bày thuật toán. Mục Kiểm tra tính đúng đắn
trả lời RQ1. Mục Phương pháp đo mô tả cách đo, và mục Kết quả và phân tích trả lời RQ2--RQ6. Báo cáo kết
thúc bằng mục Các yếu tố đe dọa độ tin cậy, mục Kết luận và phụ lục chứa bảng số liệu của mọi cấu hình
trên ba máy.
